% --- CORRECCION: Forzar interpretacion raw de la entrada ---
\UseRawInputEncoding
\documentclass[aspectratio=169, 10pt, xcolor=table]{beamer}

% --- Configuracion del Tema ---
\usetheme{Madrid}
\usecolortheme{dolphin}

% --- Paquetes necesarios ---
\usepackage[utf8]{inputenc}
\usepackage[spanish]{babel}
\usepackage{amsmath, amssymb}
\usepackage{booktabs}
\usepackage{graphicx}
\usepackage{listings}
\usepackage{tikz}
\usetikzlibrary{arrows.meta, positioning, shapes.geometric}
\usepackage{etoolbox}

% --- Ruta de Imagenes ---
\graphicspath{{./imagenes/}}

% --- Estilo para codigo Python ---
\definecolor{codegreen}{rgb}{0,0.6,0}
\definecolor{codegray}{rgb}{0.5,0.5,0.5}
\definecolor{codepurple}{rgb}{0.58,0,0.82}
\definecolor{backcolour}{rgb}{0.95,0.95,0.92}
\lstset{
    language=Python,
    backgroundcolor=\color{backcolour},
    commentstyle=\color{codegreen},
    keywordstyle=\color{blue},
    numberstyle=\tiny\color{codegray},
    stringstyle=\color{codepurple},
    basicstyle=\ttfamily\scriptsize,
    breaklines=true,
    frame=single,
    showstringspaces=false,
    literate={á}{{\'a}}1 {é}{{\'e}}1 {í}{{\'i}}1 {ó}{{\'o}}1 {ú}{{\'u}}1 {ñ}{{\~n}}1 {ü}{{\"u}}1
}

% --- Pie de pagina personalizado ---
\makeatletter
\setbeamertemplate{footline}{
  \leavevmode%
  \hbox{%
  \begin{beamercolorbox}[wd=.333333\paperwidth,ht=2.25ex,dp=1ex,center]{author in head/foot}%
    \usebeamerfont{author in head/foot}\insertshortauthor
  \end{beamercolorbox}%
  \begin{beamercolorbox}[wd=.333333\paperwidth,ht=2.25ex,dp=1ex,center]{title in head/foot}%
    \usebeamerfont{title in head/foot}Sesion 10
  \end{beamercolorbox}%
  \begin{beamercolorbox}[wd=.333333\paperwidth,ht=2.25ex,dp=1ex,right]{date in head/foot}%
    \usebeamerfont{date in head/foot}CALCULO Y ALGEBRA LINEAL \hfill \insertframenumber/\inserttotalframenumber \hspace*{0.3cm}
  \end{beamercolorbox}}%
  \vskip0pt%
}
\makeatother

% --- Datos de la portada ---
\title[SESION 10 CALCULO]{\textbf{Sesion 10}}
\subtitle{Espacios Vectoriales, Cambio de Base y Transformaciones Lineales}
\author[Prof. C. Sanchez]{Carlos Cesar Sanchez Coronel}
\institute{\textbf{CALCULO Y ALGEBRA LINEAL PARA IA}}
\date{}

% --- Eliminar barra de navegacion ---
\setbeamertemplate{navigation symbols}{}

% --- Indice al inicio de cada seccion ---
\AtBeginSection[]{
  \begin{frame}
    \frametitle{Contenido}
    \tableofcontents[currentsection]
  \end{frame}
}

\begin{document}

% --- Portada ---
\begin{frame}
    \titlepage
\end{frame}

% --- Indice general ---
\begin{frame}{Contenido}
    \tableofcontents
\end{frame}

% ============================================================================
% SECCION 1: ESPACIOS VECTORIALES
% ============================================================================
\section{Espacios Vectoriales: El Hogar de los Datos}

\begin{frame}{Definicion intuitiva}
    \begin{itemize}
        \item Un espacio vectorial es un conjunto de objetos (vectores) que se pueden sumar y multiplicar por escalares.
        \item Ejemplos: $\mathbb{R}^2$ (plano), $\mathbb{R}^3$ (espacio), matrices, funciones.
        \item En IA, los datos son vectores en un espacio de características.
    \end{itemize}
\end{frame}

\begin{frame}{Definicion formal}
    \begin{block}{Espacio vectorial sobre $\mathbb{R}$}
        Un conjunto $V$ con dos operaciones:
        \begin{itemize}
            \item Suma $+ : V\times V \to V$ (asociativa, conmutativa, neutro, inverso).
            \item Producto por escalar $\cdot : \mathbb{R}\times V \to V$ (distributivo, asociativo, identidad).
        \end{itemize}
    \end{block}
    \begin{exampleblock}{Ejemplo clasico}
        $\mathbb{R}^n$ con suma y producto por escalar usuales.
    \end{exampleblock}
\end{frame}

\begin{frame}{Subespacios}
    \begin{itemize}
        \item $W\subseteq V$ es subespacio si es cerrado bajo suma y producto por escalar.
        \item Ejemplos en $\mathbb{R}^3$: rectas y planos que pasan por el origen.
        \item $W = \{(x,y,z): x+y-z=0\}$ es un subespacio (plano).
    \end{itemize}
\end{frame}

% ============================================================================
% SECCION 2: COMBINACION LINEAL, BASE, DIMENSION
% ============================================================================
\section{Combinacion Lineal, Independencia, Base y Dimension}

\begin{frame}{Combinacion lineal}
    Dados vectores $\mathbf{v}_1,\dots,\mathbf{v}_k$ y escalares $c_1,\dots,c_k$,
    \[
    c_1\mathbf{v}_1 + \cdots + c_k\mathbf{v}_k
    \]
    es una combinacion lineal.
    \begin{exampleblock}{Ejemplo}
        $(5,8) = 2(1,2) + 1(3,4)$ es combinacion lineal de $(1,2)$ y $(3,4)$.
    \end{exampleblock}
\end{frame}

\begin{frame}{Independencia lineal}
    \begin{block}{Definicion}
        $\mathbf{v}_1,\dots,\mathbf{v}_k$ son linealmente independientes si
        \[
        c_1\mathbf{v}_1+\cdots+c_k\mathbf{v}_k = \mathbf{0} \Rightarrow c_1=\cdots=c_k=0.
        \]
    \end{block}
    \begin{exampleblock}{Ejemplo en $\mathbb{R}^3$}
        $(1,0,1)$, $(0,1,1)$, $(1,1,0)$ son independientes (determinante $-2\neq0$).
    \end{exampleblock}
\end{frame}

\begin{frame}{Base y dimension}
    \begin{block}{Base}
        Un conjunto de vectores que es linealmente independiente y genera todo el espacio.
    \end{block}
    \begin{itemize}
        \item La base canónica de $\mathbb{R}^3$: $(1,0,0)$, $(0,1,0)$, $(0,0,1)$.
        \item \textbf{Dimension}: numero de elementos de una base.
        \item $\dim(\mathbb{R}^n)=n$.
    \end{itemize}
\end{frame}

\begin{frame}{Coordenadas de un vector en una base}
    Si $\mathcal{B}=\{\mathbf{v}_1,\dots,\mathbf{v}_n\}$ es base, cualquier vector $\mathbf{v}$ se escribe de forma unica:
    \[
    \mathbf{v} = \alpha_1\mathbf{v}_1 + \cdots + \alpha_n\mathbf{v}_n
    \]
    Los $\alpha_i$ son las coordenadas de $\mathbf{v}$ en base $\mathcal{B}$, denotado $[\mathbf{v}]_{\mathcal{B}}$.
\end{frame}

% ============================================================================
% SECCION 3: CAMBIO DE BASE
% ============================================================================
\section{Cambio de Base}

\begin{frame}{Matriz de cambio de base}
    \begin{block}{De base $\mathcal{B}$ a base $\mathcal{C}$}
        Sea $\mathcal{B}=\{\mathbf{u}_1,\dots,\mathbf{u}_n\}$, $\mathcal{C}=\{\mathbf{w}_1,\dots,\mathbf{w}_n\}$. La matriz $\mathbf{P}_{\mathcal{C}\leftarrow\mathcal{B}}$ tiene como columnas $[\mathbf{u}_j]_{\mathcal{C}}$. Entonces
        \[
        [\mathbf{v}]_{\mathcal{C}} = \mathbf{P}_{\mathcal{C}\leftarrow\mathcal{B}} [\mathbf{v}]_{\mathcal{B}}.
        \]
    \end{block}
\end{frame}

\begin{frame}{Ejemplo en $\mathbb{R}^2$}
    \begin{itemize}
        \item Base canónica $\mathcal{C} = \{(1,0),(0,1)\}$.
        \item Otra base $\mathcal{B} = \{(1,1),(1,-1)\}$.
        \item $[(1,1)]_{\mathcal{C}} = (1,1)^T$, $[(1,-1)]_{\mathcal{C}} = (1,-1)^T$.
        \item $\mathbf{P}_{\mathcal{C}\leftarrow\mathcal{B}} = \begin{pmatrix}1&1\\1&-1\end{pmatrix}$.
        \item Si $[\mathbf{v}]_{\mathcal{B}} = (2,3)^T$, entonces $[\mathbf{v}]_{\mathcal{C}} = \mathbf{P}_{\mathcal{C}\leftarrow\mathcal{B}}(2,3)^T = (5,-1)^T$.
    \end{itemize}
\end{frame}

\begin{frame}{Aplicacion en NLP: embeddings}
    \begin{itemize}
        \item Las palabras se representan como vectores (embeddings).
        \item La matriz de embedding $E$ transforma un vector one-hot (base de palabras) a un espacio semántico.
        \item Esto equivale a un cambio de base: de la base de palabras a la base de características semánticas.
        \item Las operaciones algebraicas (rey - hombre + mujer $\approx$ reina) tienen sentido por las relaciones lineales en ese espacio.
    \end{itemize}
\end{frame}

% ============================================================================
% SECCION 4: TRANSFORMACIONES LINEALES
% ============================================================================
\section{Transformaciones Lineales}

\begin{frame}{Definicion}
    Una funcion $T:V\to W$ entre espacios vectoriales es lineal si:
    \begin{enumerate}
        \item $T(\mathbf{u}+\mathbf{v}) = T(\mathbf{u}) + T(\mathbf{v})$.
        \item $T(c\mathbf{v}) = cT(\mathbf{v})$.
    \end{enumerate}
    \begin{alertblock}{Importante}
        Siempre $T(\mathbf{0}) = \mathbf{0}$.
    \end{alertblock}
\end{frame}

\begin{frame}{Representacion matricial}
    Fijadas bases $\mathcal{B}_V$ y $\mathcal{B}_W$, existe una matriz $\mathbf{A}$ tal que
    \[
    [T(\mathbf{v})]_{\mathcal{B}_W} = \mathbf{A} [\mathbf{v}]_{\mathcal{B}_V}.
    \]
    \begin{itemize}
        \item La $j$-esima columna de $\mathbf{A}$ es $[T(\mathbf{v}_j)]_{\mathcal{B}_W}$ donde $\mathbf{v}_j$ es el $j$-esimo vector de la base de $V$.
    \end{itemize}
\end{frame}

\begin{frame}{Ejemplos geometricos en 2D}
    \begin{columns}[t]
        \column{0.5\textwidth}
        \textbf{Rotacion} ($\theta$):
        \[
        \begin{pmatrix}\cos\theta & -\sin\theta\\ \sin\theta & \cos\theta\end{pmatrix}
        \]
        \textbf{Escalado}:
        \[
        \begin{pmatrix}s_x & 0\\ 0 & s_y\end{pmatrix}
        \]
        \column{0.5\textwidth}
        \textbf{Cizallamiento horizontal}:
        \[
        \begin{pmatrix}1 & k\\ 0 & 1\end{pmatrix}
        \]
        \textbf{Reflexion eje y}:
        \[
        \begin{pmatrix}1 & 0\\ 0 & -1\end{pmatrix}
        \]
    \end{columns}
\end{frame}

\begin{frame}{Composicion de transformaciones}
    La composicion $S\circ T$ corresponde al producto de sus matrices (en el orden correcto):
    \[
    [S\circ T] = [S]\,[T].
    \]
    \begin{exampleblock}{Ejemplo}
        Rotar $30^\circ$ y luego escalar $(2,1)$:
        \[
        \begin{pmatrix}2&0\\0&1\end{pmatrix}
        \begin{pmatrix}\cos30^\circ&-\sin30^\circ\\ \sin30^\circ&\cos30^\circ\end{pmatrix}
        \]
    \end{exampleblock}
\end{frame}

% ============================================================================
% SECCION 5: APLICACIONES EN IA
% ============================================================================
\section{Aplicaciones en Inteligencia Artificial}

\begin{frame}{Deep Learning: capas fully connected}
    \begin{itemize}
        \item Una capa lineal: $\mathbf{y} = \mathbf{W}\mathbf{x} + \mathbf{b}$.
        \item $\mathbf{W}$ es una transformacion lineal del espacio de entrada.
        \item Las columnas de $\mathbf{W}$ son las imagenes de los vectores base.
    \end{itemize}
    \begin{center}
        \begin{tikzpicture}[scale=0.8]
            \draw[->] (-1,0) -- (2,0) node[right] {$x_1$};
            \draw[->] (0,-1) -- (0,2) node[above] {$x_2$};
            \draw[thick, blue] (0,0) -- (1,0) node[below] {$\mathbf{e}_1$};
            \draw[thick, blue] (0,0) -- (0,1) node[left] {$\mathbf{e}_2$};
            \draw[->, red] (3,1) -- (5,2) node[midway, above] {$\mathbf{W}$};
            \draw[->] (6,-1) -- (9,2);
            \draw[->] (6,2) -- (9,1);
        \end{tikzpicture}
    \end{center}
\end{frame}

\begin{frame}{NLP: embeddings y cambio de base}
    \begin{itemize}
        \item Matriz de embedding $E$: cada columna es el embedding de una palabra.
        \item Operaciones como $\mathbf{rey} - \mathbf{hombre} + \mathbf{mujer} \approx \mathbf{reina}$ son transformaciones lineales en el espacio semántico.
        \item Un cambio de base corresponde a una reparametrizacion del espacio.
    \end{itemize}
\end{frame}

\begin{frame}{Vision por computador: data augmentation}
    \begin{itemize}
        \item Para aumentar datos, aplicamos transformaciones afines a imagenes: rotacion, escalado, cizallamiento.
        \item Cada transformacion es lineal en coordenadas homogeneas.
        \item Ejemplo: rotar $45^\circ$ y trasladar $(10,0)$:
        \[
        \begin{pmatrix}\cos45^\circ&-\sin45^\circ&10\\ \sin45^\circ&\cos45^\circ&0\\0&0&1\end{pmatrix}
        \]
    \end{itemize}
\end{frame}

% ============================================================================
% SECCION 6: VISUALIZACION CON PYTHON
% ============================================================================
\section{Visualizacion con Python}

\begin{frame}[fragile]{Transformaciones en 2D con Matplotlib}
    \begin{lstlisting}
import numpy as np
import matplotlib.pyplot as plt

def plot_transform(A, puntos, titulo):
    puntos_trans = A @ puntos
    plt.figure(figsize=(8,4))
    plt.subplot(1,2,1)
    plt.scatter(puntos[0,:], puntos[1,:], c='b', alpha=0.5)
    plt.quiver(0,0,1,0, angles='xy', scale_units='xy', scale=1, color='r', label='e1')
    plt.quiver(0,0,0,1, angles='xy', scale_units='xy', scale=1, color='g', label='e2')
    plt.xlim(-3,3); plt.ylim(-3,3); plt.grid(); plt.axis('equal')
    plt.title('Original')
    plt.subplot(1,2,2)
    plt.scatter(puntos_trans[0,:], puntos_trans[1,:], c='b', alpha=0.5)
    e1_t = A @ [1,0]
    e2_t = A @ [0,1]
    plt.quiver(0,0, e1_t[0], e1_t[1], angles='xy', scale_units='xy', scale=1, color='r')
    plt.quiver(0,0, e2_t[0], e2_t[1], angles='xy', scale_units='xy', scale=1, color='g')
    plt.xlim(-3,3); plt.ylim(-3,3); plt.grid(); plt.axis('equal')
    plt.title('Transformado')
    plt.suptitle(titulo)
    plt.tight_layout()
    plt.show()
    \end{lstlisting}
\end{frame}

\begin{frame}[fragile]{Ejemplos de transformaciones}
    \begin{lstlisting}
# Crear puntos de un cuadrado
x = np.linspace(-1,1,10)
y = np.linspace(-1,1,10)
X, Y = np.meshgrid(x,y)
puntos = np.vstack([X.ravel(), Y.ravel()])

# Rotacion 30 grados
theta = np.radians(30)
A_rot = np.array([[np.cos(theta), -np.sin(theta)],
                  [np.sin(theta), np.cos(theta)]])
plot_transform(A_rot, puntos, 'Rotación 30°')

# Escalado
A_scale = np.array([[2,0],[0,0.5]])
plot_transform(A_scale, puntos, 'Escalado (2x, 0.5y)')

# Cizallamiento horizontal
A_shear = np.array([[1,1.5],[0,1]])
plot_transform(A_shear, puntos, 'Cizallamiento horizontal')
    \end{lstlisting}
\end{frame}

\begin{frame}[fragile]{Capa lineal con NumPy}
    \begin{lstlisting}
class LinearLayer:
    def __init__(self, in_features, out_features):
        self.W = np.random.randn(out_features, in_features) * 0.01
        self.b = np.zeros((out_features, 1))

    def forward(self, X):
        # X: (in_features, batch_size)
        return self.W @ X + self.b

# Ejemplo: batch de 5 vectores de 3 dims, capa de salida 2 dims
batch = np.random.randn(3,5)
layer = LinearLayer(3,2)
salida = layer.forward(batch)
print("Entrada shape:", batch.shape)
print("Salida shape:", salida.shape)
    \end{lstlisting}
\end{frame}

\begin{frame}[fragile]{Cambio de base con NumPy}
    \begin{lstlisting}
# Base B: columnas (1,1) y (1,-1)
P = np.array([[1, 1],
              [1, -1]])

# Vector en base canonica
v_canon = np.array([5, -1])

# Coordenadas en base B resolviendo P x = v_canon
v_B = np.linalg.solve(P, v_canon)
print("Coordenadas en base B:", v_B)   # [2, 3]

# Verificacion
print("Reconstruccion:", P @ v_B)      # [5, -1]
    \end{lstlisting}
\end{frame}

% ============================================================================
% SECCION 7: NOTACION EN PAPERS
% ============================================================================
\section{Notacion en Papers Cientificos}

\begin{frame}{Notacion comun}
    \begin{itemize}
        \item $\mathbf{h} = \mathbf{W}\mathbf{x} + \mathbf{b}$: capa lineal (fully connected).
        \item $\mathbf{W} \in \mathbb{R}^{d_{\text{out}}\times d_{\text{in}}}$ matriz de pesos.
        \item $\mathbf{E}_{\text{word}}$: matriz de embedding.
        \item $\mathbf{T}_\theta(\mathbf{p}) = \mathbf{R}_\theta \mathbf{p} + \mathbf{t}$: transformacion afín.
        \item Descomposicion SVD: $\mathbf{W} = \mathbf{U}\mathbf{D}\mathbf{V}^\top$ (cambio de base ortonormal).
    \end{itemize}
\end{frame}

% ============================================================================
% SECCION 8: EJERCICIOS PROPUESTOS
% ============================================================================
\section{Ejercicios Propuestos}

\begin{frame}{Ejercicios}
    \begin{enumerate}
        \item Demuestra que el conjunto de matrices $2\times2$ es un espacio vectorial. ¿Cual es su dimension?
        \item Determina si $(1,2,3)$, $(4,5,6)$, $(7,8,9)$ son linealmente independientes.
        \item Encuentra la matriz de la transformacion lineal que refleja un vector respecto a la recta $y=x$.
        \item Genera 100 puntos aleatorios en $[-1,1]\times[-1,1]$ y aplicales una rotacion de $60^\circ$, un escalado $(2,0.5)$ y un cizallamiento $0.8$. Visualiza.
        \item Investiga como se representa una convolucion como transformacion lineal (matriz de Toeplitz).
    \end{enumerate}
\end{frame}

% ============================================================================
% SECCION 9: RESUMEN
% ============================================================================
\section{Resumen}

\begin{frame}{Lo que hemos aprendido}
    \begin{exampleblock}{Conceptos clave}
        \begin{itemize}
            \item \textbf{Espacios vectoriales}: estructura matematica de los datos.
            \item \textbf{Base y dimension}: coordenadas unicas.
            \item \textbf{Cambio de base}: matriz $\mathbf{P}_{\mathcal{C}\leftarrow\mathcal{B}}$.
            \item \textbf{Transformaciones lineales}: representacion matricial, ejemplos geometricos.
            \item \textbf{Aplicaciones}: capas fully connected, embeddings, data augmentation.
            \item \textbf{Visualizacion}: Python con NumPy y Matplotlib.
        \end{itemize}
    \end{exampleblock}
\end{frame}

\begin{frame}{Conexion con futuros cursos}
    \begin{itemize}
        \item \textbf{Deep Learning}: redes profundas apilan transformaciones lineales.
        \item \textbf{NLP}: espacios de embeddings, analogias.
        \item \textbf{Vision}: transformaciones afines y aumento de datos.
        \item \textbf{Algebra lineal avanzada}: descomposiciones espectrales, PCA.
    \end{itemize}
\end{frame}

% --- Diapositiva final ---
\begin{frame}
    \centering
    \Huge \textbf{Gracias!}

\end{frame}

\end{document}